% from ega4-12-fr.tex
\oldpage[IV-3]{173}
\section{Étude des fibres des morphismes plats de présentation finie}
\label{section:IV.12-fr}

\setcounter{subsection}{-1}
\subsection{Introduction}
\label{subsection:IV.12.0-fr}

\emph{Nous utiliserons dans tout ce paragraphe les notations générales de}
(9.4.1).

\begin{env}[12.0.1]
\label{IV.12.0.1-fr}
Étant donné un morphisme localement de présentation finie $f:X\to Y$, nous
avons vu (9.9) que pour certaines propriétés locales $\mathbf P$ de préschémas
sur des corps ou de Modules sur ces préschémas, l'ensemble $E$ des $x\in X$
tels que la propriété $\mathbf P$ ait lieu, pour la fibre $X_{f(x)}$, au point
$x$ de cette fibre, est localement constructible dans $X$. Nous nous proposons
de montrer que, pour la plupart de ces propriétés, si l'on suppose de plus que
le morphisme $f$ est plat, alors l'ensemble $E$ est même ouvert dans $X$. De
même ((9.2) à (9.8)) nous avons montré que si $f$ est de présentation finie,
et si cette fois $\mathbf P$ désigne certaines propriétés globales de
préschémas sur des corps ou de Modules sur ces préschémas, l'ensemble $F$ des
$y\in Y$ tels que la propriété $\mathbf P$ ait lieu pour la fibre $X_y$ est
localement constructible dans $Y$. Nous montrerons que si l'on suppose de plus
que le morphisme $f$ est propre et plat, $F$ est même ouvert dans $Y$.
\end{env}

\begin{env}[12.0.2]
\label{IV.12.0.2-fr}
La méthode générale de démonstration des propriétés en question comporte trois
étapes. On se ramène d'abord au cas où $Y$ est affine et $X$ de présentation
finie~; puis~:
\begin{enumerate}[label=\emph{\Alph*)}]
  \item À l'aide de (8.9.1) (et éventuellement d'autres résultats du \S~8) et
  de (11.2.6), on se ramène au cas où $X$ et $Y$ sont noethériens.
  \item On applique les résultats du \S~9 rappelés en (12.0.1) prouvant que
  $E$ (resp. $F$) est constructible.
  \item Pour voir que $E$ est ouvert, il suffit, en vertu de
  ($\mathbf 0_{\mathrm{III}}$, 9.2.5) de montrer que si $x\in E$, alors toute
  générisation $x'$ de $x$ appartient aussi à $E$. Utilisant ($\mathbf{II}$,
  7.1.7), on voit, puisque $Y$ est noethérien, qu'il existe un spectre d'anneau
  de valuation discrète $Y_1$ et un morphisme $h:Y_1\to X$ tel que, si $y_1$
  (resp. $y'_1$) est le point fermé (resp. le point générique) de $Y_1$, on ait
  $h(y_1)=x$, $h(y'_1)=x'$. On fait alors le changement de base
  $g=f\circ h:Y_1\to Y$~; compte tenu des résultats des \S\S~4 et 6 sur les
  préschémas localement noethériens sur des corps et les changements du corps
  de base, on est ramené à prouver l'assertion en question pour un point $x_1$
  de $X_1=X\times_YY_1$ au-dessus de $y_1$ et pour une générisation $x'_1$ de
  $x_1$ au-dessus de $y'_1$. Comme $Y_1=\operatorname{Spec}(A)$, où $A$ est
  un anneau de valuation discrète, d'uniformisante $t$, on doit en définitive,
  pour un $A$-module $M$, prouver une propriété de $M$, sachant que $M/tM$ a
  la même propriété et que $t$ est $M$-régulier (ce qui résulte de l'hypothèse
  de platitude)~; on utilise pour cela les résultats de (3.4) et de (5.12). On
  procède de même pour l'ensemble $F\subset Y$.
\end{enumerate}
\end{env}

